% Transcription — fonds Alexandre Grothendieck, shelfmark 161-4, batch 2 (pages 21-31).
% Established from the facsimile archives/batches/161-4/batch-02.pdf (a mirror of
% https://grothendieck.umontpellier.fr/161-4.pdf), per the skill
% transcribe-grothendieck. The batch file carries no cover sheet: PDF page N is
% archivists' page 20 + N, checked against the pencilled numbers.
%
% Pass: Opus 5.5 (claude-opus-5-5), 2026-09-22 — first pass, unchecked against the pages by a human.
%
% What the batch holds, in the archivists' order. All eleven pages are
% mathematics; none is skipped.
%   21     Black ink, squared paper: « Topologie des spectres » — noetherian
%          spaces, Spec functoriality (φ(x_p), φ^{-1}(V(J)), φ^{-1}(Y_f)),
%          closed and open immersions, Spec of a product, and the warning
%          that |Spec(X ×_S Y)| → |Spec X| ×_{|Spec S|} |Spec Y| is onto but
%          not injective (fibre = Spec k(x) ⊗_{k(s)} k(y)), with four sketches.
%   22-24  Black ink, spiral-pad sheets: the start of the course proper.
%          Page 22 opens in English (« Prerequisites for the courses I and
%          II » — I: Introduction to A.G., II: algebraic groups; a summer
%          course), then numbered items (1)-(8) in French: P_I = k[(T_i)],
%          k^I ≃ Hom_{k-alg}(P_I, k), arithmetic vs geometric point of view,
%          the functor E^I_k : k' ↦ k'^I, V_S ⊂ E^I_k, closed subspaces =
%          affine algebraic spaces, A_J = P_I/J, representable functors V_A,
%          anti-equivalence with k-algebras (5), restriction to fields and
%          radical ideals J~ (6), the Nullstellensatz (6'), A as the ring of
%          functions Hom(V_A, E^1_k) (8).
%   25     A pencil sheet crossed through by long diagonals, on another
%          subject (topos: Hom_X(F, L_X), X/F, I^E, « Topos associé à
%          l'espace compact I^E »). Read only in fragments.
%   26-28  Blue ink, plain paper: items (9)-(11) — projective limits in
%          Aff_k ↔ inductive limits in Alg_k (products, fibre products, base
%          change, kernels, cokernels of double arrows, filtrant limits),
%          stabilisers, Sl(n) = Ker det; inductive limits in Aff_k (G/B = e_k
%          « ce qui est idiot ! »), finite sums and « connected » k', the
%          empty space as strict initial object; (11) the spectrum: points =
%          prime ideals.
%   29-31  Blue ink, same paper: (b) the Zariski topology on X = Spec A,
%          V(S), closure of a point, quasi-compactness, irreducible ⇔
%          integral, sobriety; the structure sheaf via M_f, the relation
%          f ≺ g, the sheaf condition, stalks A_p local.
%
% Notation fixed for the batch. His blackboard E (E^I_k, E^1_k) is
% \mathbb{E}; his gothic X for the affine algebraic space is \mathfrak{X};
% P_I, A_J, V_S, V_J, V_A as written; J~ (radical) is \widetilde{J}; ideals
% p are \mathfrak{p}; « lim » over a left arrow is \varprojlim, over a right
% arrow \varinjlim; Aff_k, Alg_k are \mathrm{Aff}_k, \mathrm{Alg}_k; the
% structure sheaf, underlined O, is \underline{\mathcal{O}}_X; the preorder
% he writes f ≺ g is \prec. Item numbers circled on the page are set bold
% in parentheses.
%
% Slips left as written, flagged: « algèbres sur A » for « sur k » (p. 23),
% « qui contiennent A » for « S » (p. 29), « f_i ∈ I » (p. 30).
%
% Doubtful, not settled: the margins of pages 24 and 31, the whole of the
% pencil page 25, several struck words throughout, the connective phrase of
% page 28's first lines.
%
% For the modernised reading. Page 22 shows the notes to be course preparation
% for a two-part summer course, in English, « I) Introduction to A.G. »
% then « II) Introduction to algebraic groups »; the French numbered items
% (1)-(11) on pages 22-28 are its outline, and pages 21 and 29-31 carry the
% spectrum's topology and structure sheaf. Page 30's margin (« 10')
% Fonctorialité … 11) ») shows he meant to renumber. Page 21 continues the
% spectrum material and is likely later in the argument than page 22 despite
% the archivists' order; page 25 does not belong to the course.
\documentclass[11pt,a4paper]{article}
% Le préambule commun à toutes les transcriptions du fonds Grothendieck.
%
% Il tient en une idée : **l'incertitude se compose, elle ne se résout pas.**
% Une transcription qui lisse un mot douteux a détruit la seule chose pour
% laquelle on la faisait — un lecteur doit pouvoir savoir, à chaque endroit,
% ce qui était sur le feuillet et ce qui a été ajouté. Les six macros qui
% suivent sont donc l'appareil critique, et non de la décoration.
%
% Les mêmes commandes sont reconnues par `scripts/render.mjs`, qui produit la
% vue de lecture du site. Une macro ajoutée ici doit l'être aussi là-bas,
% faute de quoi le rendu échouera bruyamment — ce qui est voulu : un rendu
% silencieusement faux est pire qu'un rendu absent.

\ProvidesPackage{grothendieck}[2026/08/08 Transcriptions du fonds Grothendieck]

\usepackage{amsmath,amssymb,amsthm}
\usepackage{mathrsfs}
\usepackage{stmaryrd}
% Les diagrammes commutatifs sont partout dans ce fonds : c'est la langue
% même de la géométrie algébrique de Grothendieck, et une transcription qui
% les rendrait en prose ne transcrirait rien.
\usepackage{tikz-cd}
% Français par défaut — c'est la langue des feuillets — mais l'anglais est
% chargé aussi : la traduction et le résumé sont des documents anglais, et la
% typographie française (espaces avant les deux-points) y serait une faute.
% Ces documents basculent par \selectlanguage{english} en tête de corps.
\usepackage[english,french]{babel}
\usepackage{fontspec}
\usepackage{geometry}
\geometry{a4paper,margin=25mm,marginparwidth=28mm}
\usepackage{marginnote}
\usepackage{xcolor}
% ulem, not soul. `soul`'s \st and \ul are built on a character-by-character
% scan that XeLaTeX's font machinery breaks: the document fails with an
% undefined control sequence rather than degrading. ulem does the same two jobs
% and is Unicode-engine safe. `normalem` stops it hijacking \emph, which the
% transcriptions use for Grothendieck's own emphasis.
\usepackage[normalem]{ulem}
\usepackage{hyperref}
\hypersetup{colorlinks=true,linkcolor=black,urlcolor=black,pdfborder={0 0 0}}

% --- Métadonnées du lot -----------------------------------------------------
% Reprises telles quelles par le script de rendu, qui les affiche en tête de
% la vue de lecture. Elles fixent ce que le fichier prétend couvrir : sans
% elles, une transcription flotte sans point d'attache dans un fonds de seize
% mille feuillets.

\newcommand{\folder}[1]{\gdef\grothFolder{#1}}
\newcommand{\batch}[1]{\gdef\grothBatch{#1}}
\newcommand{\leaves}[2]{\gdef\grothFirst{#1}\gdef\grothLast{#2}}
\newcommand{\foldertitle}[1]{\gdef\grothFoldertitle{#1}}
\gdef\grothFolder{?}\gdef\grothBatch{?}\gdef\grothFirst{?}\gdef\grothLast{?}\gdef\grothFoldertitle{}

% --- Le résumé --------------------------------------------------------------
%
% Il ouvre la lecture modernisée, sous le titre « Résumé », et il est composé
% en petit. Ce n'est pas un ornement : c'est le seul passage du document qui
% ne soit pas des mathématiques, et un lecteur doit pouvoir le reconnaître
% avant de l'avoir lu — puis le sauter s'il connaît déjà le sujet, ou s'y
% arrêter s'il ne le connaît pas. Le filet qui le suit marque la reprise du
% corps.

\newenvironment{resume}
  {\small\setlength{\parskip}{0.45em}}
  {\par\medskip\hrule\medskip}

% --- Mots-clés ---------------------------------------------------------------
%
% En anglais, en clôture du résumé de la lecture modernisée : le vocabulaire
% moderne sous lequel chercher ce que les feuillets construisent. Le manifeste
% du site (`npm run manifest`) les extrait pour étiqueter la cote entière —
% cette ligne est la source unique des tags, il n'y a pas de fichier à côté.
\newcommand{\keywords}[1]{\par\smallskip\noindent
  {\footnotesize\textsc{Keywords} --- \textit{#1}}\par}

% --- L'appareil critique ----------------------------------------------------

% Le numéro de feuillet, en marge. C'est l'ancre qui permet de régler un
% désaccord sur une formule en regardant *un* feuillet plutôt qu'une chemise
% de six cents. Le site s'en sert aussi pour tourner les pages du fac-similé
% quand on fait défiler la transcription.
\newcommand{\leaf}[1]{%
  \marginnote{\footnotesize\color{gray!70}\textsf{#1}}%
  \label{leaf:#1}\ignorespaces}

% Illisible : on ne devine pas. Un mot inventé qui se lit comme les autres est
% le pire résultat possible.
\newcommand{\ill}{\textcolor{red!60!black}{[\,\ldots\,]}}

% Lecture douteuse : le mot est donné, et signalé comme incertain.
\newcommand{\uncertain}[1]{\uline{#1}}

% Ajout de l'éditeur — une lettre manquante, un mot sous-entendu.
\newcommand{\add}[1]{\textcolor{blue!50!black}{[#1]}}

% Biffé par Grothendieck lui-même. Ce qu'il a rayé fait partie du document :
% une rature dit souvent où la pensée a changé de direction.
\newcommand{\struck}[1]{\sout{#1}}

% Note du transcripteur, sur l'état matériel du feuillet.
\newcommand{\note}[1]{\par\smallskip{\small\color{gray!60!black}\textit{[#1]}}\par\smallskip}

% Note marginale de Grothendieck, distincte du corps du texte.
\newcommand{\marginal}[1]{\par\smallskip\noindent\hspace{1em}%
  {\small\color{gray!60!black}#1}\par\smallskip}

% --- Titre ------------------------------------------------------------------

\AtBeginDocument{%
  \begin{center}
    {\large\bfseries Fonds Alexandre Grothendieck --- cote n\textsuperscript{o}~\grothFolder}\\[2pt]
    {\itshape\grothFoldertitle}\\[4pt]
    {\small feuillets \grothFirst--\grothLast\ (lot \grothBatch)}\\[2pt]
    {\footnotesize\color{gray!60!black}Transcription établie d'après le fac-similé numérisé
     par l'Université de Montpellier.\\
     Les lectures incertaines sont soulignées, les illisibles notées [\,\ldots\,],
     les ajouts entre crochets.}
  \end{center}
  \vspace{1em}}

\endinput


\folder{161-4}
\batch{2}
\pages{21}{31}
\dating{[après 1961] — le groupe « Dossiers rassemblés par Grothendieck sur différentes thématiques » (161-1 à 162-6) est daté [après 1961-vers 1977]}
\watermark{Édition de démonstration}
\foldertitle{Introduction à géométrie algébrique : notes manuscrites (s.d.).}

\begin{document}

\page{21}
\emph{Topologie des spectres}

Remarque : espace noethérien $=$ \struck{tout} espace quasicompact tel que a) tout sous-espace fermé ait un nb fini de comp. irréductibles b) toute suite décroissante de sous-espaces fermés \uncertain{est stationnaire}.

Composantes connexes d'un espace top. n'ayant qu'un nb fini de composantes irréductibles (p.ex. noeth.) --- plus gén. où \uncertain{l'ens.} des comp. irréd. est loc. fini.

\struck{Topologie \ill{} $A_f$ \ill{} $X_f$. L'espace topologique $\operatorname{Spec} A$ pour $A$ \uncertain{noethérien}.}

base de topologie formée des $X_f$ \struck{/ homéomorphismes $X_f \simeq \operatorname{Spec} A_f$}

L'espace topologique $\operatorname{Spec} A$ pour $A$ \uncertain{noethérien} \ill{}

Cf. $u : A \to B$, $Y = \operatorname{Spec} A \xleftarrow{\ \varphi\ } X = \operatorname{Spec} B$
\[
\begin{cases}
\varphi(x_{\mathfrak{p}}) = y_{\varphi^{-1}(\mathfrak{p})} \\
\varphi^{-1}(V(J)) = V(u(J)B) \\
\varphi^{-1}(Y_f) = X_{u(f)}
\end{cases}
\]
\note{dans la dernière ligne l'indice est surchargé : $\varphi(f)$ corrigé, semble-t-il, en $u(f)$ ; même surcharge sous le $V(J)$ de la deuxième ligne}

\emph{Ex} a) $u : A \to B$ surjectif, $\varphi$ imm. fermée.

b) $u : A \to A_f$ \qquad $\varphi$ immersion ouverte \qquad $\operatorname{Spec}(A_f) \simeq X_f$

dans les deux cas, les \uncertain{ext.} résiduelles sont triviales, cas $A \to B$ un épim. d'anneaux [i.e. $\mathfrak{X} \to \mathfrak{Y}$ un monomorphisme d'esp. alg. affines]

c) $\operatorname{Spec} \prod_{I} A_i \simeq \coprod_i \operatorname{Spec}(A_i)$, \qquad $X \simeq \dot{\coprod_i} X_i$

d) Attention, le \emph{foncteur $\mathfrak{X} \to \operatorname{Spec} \mathfrak{X}$ ne commute pas aux produits ou produits fibrés} [vrai sur un corps alg. clos]
\[
\lvert \operatorname{Spec} \mathfrak{X} \times_S \mathfrak{Y} \rvert \longrightarrow \lvert \operatorname{Spec}(\mathfrak{X}) \rvert \times_{\lvert \operatorname{Spec} S \rvert} \lvert \operatorname{Spec} \mathfrak{Y} \rvert
\]
est une surjection, mais non pas injection : la \struck{\ill{}} fibre de $x, y$ \struck{\ill{}} est \uncertain{le spectre} : $x \times_S y \simeq \operatorname{Spec} k(x) \otimes_{k(s)} k(y)$ [qui est $\neq \emptyset$].

\begin{tikzcd}[column sep=small, row sep=small]
& \mathfrak{X} \times_S \mathfrak{Y} \arrow[dl] \arrow[dr] & \\
x \in \mathfrak{X} \arrow[dr, no head] & & \mathfrak{Y} \ni y \arrow[dl, no head] \\
& S \ni s &
\end{tikzcd}

\begin{tikzcd}[column sep=small, row sep=small]
& K' \text{ corps} \arrow[dl, no head] \arrow[dr, no head] & \\
k(x) \arrow[dr, no head] & & k(y) \arrow[dl, no head] \\
& k(s) &
\end{tikzcd}

\begin{tikzcd}[column sep=small, row sep=small]
& Z = V_{K'} \arrow[dl, no head] \arrow[dr, no head] & \\
V_{k(x)} \to V_A & & V_B \leftarrow V_{k(y)} \\
x \to X \arrow[dr] & & Y \leftarrow y \arrow[dl, no head] \\
& S = V_k &
\end{tikzcd}

\begin{tikzcd}[column sep=small, row sep=small]
& C \arrow[dl, no head] \arrow[dr, no head] & \\
A \arrow[dr, no head] & & B \arrow[dl, no head] \\
& k &
\end{tikzcd}

\marginal{donne exemple du produit $\mathbb{E}^1 \times \mathbb{E}^1$}
\note{les quatre croquis sont redessinés : à gauche, le produit $Z = V_{K'}$ de deux points $x$, $y$ au-dessus de $S = V_k$ ; au centre les algèbres correspondantes ; à droite le produit fibré et le corps composé $K'$ de $k(x)$ et $k(y)$ sur $k(s)$ ; les traits sans flèche sont ceux de la page}

\struck{Mais} (Sur un corps alg. clos, [pour esp. alg. définis de t.f.] le Spec Maximal commute aux produits finis [formation des pts ensembliste], mais sa top. n'est pas la topologie produit.

\page{22}
Prerequisites for the courses I and II

I) Introduction to A.G.\ : Category theory, basic algebra including some commutative algebra. Bourbaki is a good reference. Most facts we are using will be recalled.

II) Introduction to algebraic groups. Familiarity with the basic notions of A.G. will be assumed.

A.G. is a messy subject, \uncertain{even more foundations}, we cannot give anything like a detailed exposition of all of it in our summer course. We will have to give rather a survey of some key points, being rather sketchy for proofs --- most of which are straightforward. We can fill in whatever is required. Homework will be necessary for those who \uncertain{do} not have already a fair background.

Question is how to relate the two courses. The content of I are prerequisites for II!

\textbf{(1)} $k$
\[
P_I = k[(T_i)_{i \in I}] \qquad\qquad I, J \text{ finis ou infinis}
\]
\[
k^I \simeq \operatorname{Hom}_{k\text{-alg}}(P_I, k) \qquad \struck{\text{l'esp. affine sur } k \text{ de dim } I \ \ill{}}
\]
\[
f \in P_I, \quad x = (x_i)_{i \in I} \in k^I \qquad f(x) = \varepsilon_x(f)
\]
\[
(1) \qquad F_j(t) = 0 \quad j \in J \qquad\qquad V((F_j))
\]
\note{la première ligne s'écrit à la page $P[(T_i)_{i \in I}]$, sans doute pour $k[\ldots]$}

\textbf{(2)} (a) pt de vue arithmétique : très différent suivant la nature de $k$ : …, $\mathbb{R}$, $\mathbb{C}$

\textbf{(3)} (b) pt de vue géométrique, issu de la géométrie sur $\mathbb{R}$ et sur $\mathbb{C}$ : on \uncertain{considère} les solutions de (1) \struck{dans} toute \struck{extension} $k$-algèbre $k'$ ; et comment elles varient ; et les propriétés de (1) qui sont stables par de tels \emph{changements}…
\[
\mathbb{E}^I_k : k' \longmapsto k'^I \qquad \mathbb{E}^I_k(k') = k'^I
\]
\[
V_S(k') \subset \mathbb{E}^I_k(k') \qquad\qquad V_S \subset \mathbb{E}^I_k
\begin{cases}
\text{étude indépendante du plongement} \\
\text{étude du plongement}
\end{cases}
\]

\page{23}
\emph{Sous-espace \struck{algébrique} fermé de $\mathbb{E}^I_k$} $=$ \emph{espace alg. affine}

\textbf{(4)} $V_S = V_J$ \quad $J$ idéal engendré, c'est le syst. d'équations le plus naturel.

\marginal{référence au « th. de la base » de Hilbert ($k$ corps \struck{\ill{}} \uncertain{ou gén.}, noeth.) $\Longrightarrow k[T_1 \ldots T_n]$ noeth.}
\[
A_J = P_I / J
\]
\[
\begin{array}{ccc}
k'^I & \simeq & \operatorname{Hom}_{k\text{-alg}}(P_I, k') \\
\cup & & \cup \\
V_J(k') & \simeq & \{ u : P_I \to k' \mid u(J) = 0 \ \struck{\ill{}} \ (\text{i.e. s'annule sur les } F_j) \} \\
& & \wr\! \mid \\
& & \operatorname{Hom}_k(A_J, k')
\end{array}
\]

\struck{\ill{}} Les foncteurs de $\mathcal{C}$ \ill{} forme $V_S$ (\ill{} $S$ convenable) sont, à isom. près, les foncteurs représentables $V_A$ --- donner un isom. entre $V_A$ et un $V_S$ \struck{\ill{}} revient au même que de se donner un isom $A \simeq P_I / J$ ($J$ idéal engendré par $S$) \struck{\ill{} ; on donne un plongement \ill{} $V_S$ fermé de $\mathbb{E}_k$ ?} [Plus précisément, \ill{} plongement $V_A \hookrightarrow \mathbb{E}^I$ correspond à un épimorphisme d'anneaux $P_I \to A$ (mais un tel épimorphisme n'est pas en général surjectif…] \uncertain{identifiés} aux hom. surjectifs $P_I \to A$, on revient aux \uncertain{systèmes de générateurs} $(a_i)_{i \in I}$ de $A$ comme une \emph{algèbre sur $k$}.
\marginal{utilisant la notion de foncteurs représentables --- Rmq}
\note{le passage entre $V_S$ et le crochet est surchargé : une ligne biffée, puis la phrase « se donner un plongement » récrite au-dessus ; seul ce qui survit est lu}

\textbf{(5)} \emph{Cor} La cat. des espaces alg. affines sur $k$ est antiéquivalente à celle des algèbres sur $A$\note{sic : on attend « sur $k$ »}. Ceux qui se plongent comme sous-espaces alg. fermés dans quelque $\mathbb{E}^n_k$ ($n$ entier $\geqslant 0$ fini) \struck{sont ceux} correspondant aux $A$ \struck{qui sont des $k$-alg. de type fini} \uncertain{de type fini}.

Les \struck{\ill{}} sous-espaces algébriques de $\mathbb{E}^I_k$ correspondent aux idéaux de $P_I$. \struck{\ill{}} [$J$ connu quand on connaît $V_J$…]
\[
V_J \subset V_{J'} \Longleftrightarrow J' \subset J
\qquad \Longleftarrow \quad
\Bigl\{ J = \{ f \in P_I \mid f(x) = 0 \ \forall\, k'/k \text{ et } x \in V_J(k') \} \Bigr.
\]

\page{24}
\textbf{(6)} Ce qui se passe si on se restreint au cas où $k'$ est un corps [on est tenté de le faire du moins si $k$ est un corps]. \struck{\ill{}} Soit $V'_S = V_S \mid$ corps sur $k$
\[
V'_{S} = V'_{J} = V'_{\widetilde{J}}
\]
\note{les trois indices portaient d'abord un « corps » biffé : $V_{\text{corps}, S} = V_{\text{corps}, J} = V_{\text{corps}, \widetilde{J}}$, récrit avec l'accent}
Donc $J \mapsto V'_J$ \uncertain{n'est plus} injective. Mais sa restriction à l'ens des idéaux $J$ tels que $J = \widetilde{J}$ l'est, car
\[
\widetilde{J} = \{ f \in P_I \mid f(x) = 0 \ \forall x \in V_J(k'),\ k'/k \text{ un corps} \}
\]
[Cela \struck{signifie} équivaut à l'énoncé suivant d'AC

\emph{Prop} ($\forall J \subset P_I$ idéal) $\widetilde{J}$ est l'intersection des noyaux des hom. de $P_I / J$ dans des corps

$\Updownarrow$

$\forall$ anneau $A$, l'intersection des hom. de $A$ dans des corps [i.e. l'intersection des idéaux premiers] est \uncertain{l'ens.} des éléments nilpotents.] Donc
\[
\widetilde{J} \subset \widetilde{J}' \Longleftrightarrow V(J') \subset V(J)
\]
Mêmes conclusions (si on se borne [pour $I$ fini] : the sous-catégorie $\mathcal{C}$ de la catégorie des $k$-algèbres \uncertain{réduites} pourvu que \ill{} \ill{} finis que $k$ soit assez grand \ill{}, l'analogue de la propr. précédente reste valable

\textbf{(6')} (p. ex. $\mathcal{C}$ contient les extensions de corps,] ou \struck{\ill{}} si $I$ est fini, s'il contient les ext. finies des corps résiduels de $k$, --- p. ex. on aura : $k$ si $k$ est un corps alg. clos…)
\marginal{Nullstellensatz}
\marginal{Dans une alg. de t.f. sur un corps $k$, tt idéal premier (radical) est intersection d'idéaux maximaux \ill{} \ill{} algébriques \ill{}…}

\textbf{(7)} Importance des \struck{\ill{}} \ill{}

\textbf{(8)}
\[
A \simeq \operatorname{Hom}_{k\text{-alg}}(k[T], A) \simeq \operatorname{Hom}(V_A, \mathbb{E}^1_k)
\qquad / \qquad
A \text{ l'\emph{anneau des fonctions sur l'espace algébrique affine}} \ X.
\]
[Si $A = P_I / J$, $f \in A$ se relève en $g \in P_I$, alors $\longleftrightarrow$
\[
V_A(k') \xrightarrow{\ f_{k'}\ } k' \quad \text{est induit par} \quad \mathbb{E}^I(k') = k'^I \xrightarrow{\ g_{k'}\ } \quad \text{défini par } g ]
\]
La structure d'anneau de $A$ provient de la structure d'anneau sur le foncteur $\mathbb{E}^1_k$.
\marginal{abond. de $A$ --- définir $f \in A$ : $f(x)$ pour $x \in V_A(k')$ ; $(f+g)(x) = f(x) + g(x)$, $(fg)(x) = f(x) g(x)$ \ill{}}
\note{la marge gauche, écrite en biais et surchargée, n'est lue qu'en partie}

\page{25}
\note{feuillet au crayon, d'une autre main de travail que les pages qui l'entourent (topos, et non espaces algébriques affines) ; il est barré de longs traits diagonaux croisés, et n'est lu que par fragments ; les flèches sont redessinées}

\begin{tikzcd}[column sep=large]
X/F \arrow[d, "\wr" description, no head] & X'/F' \arrow[d, no head] \\
\underline{\operatorname{Hom}}_X(F, \mathbb{L}_X) \arrow[r] & X' \\
X \arrow[u, hook] \arrow[ur] &
\end{tikzcd}

$\underline{\operatorname{Hom}}$ \qquad \struck{\ill{}} $F$ \qquad $\mathfrak{X} / \underline{\operatorname{Hom}}(F, \mathbb{L})$ \qquad $X/F$

$E$ \uncertain{un ordinal} \ill{} $\{0 < 1\}$ \ill{} \ill{}

$E$ \ill{} \ill{} \uncertain{sommes} des $E_{X'}$ \ill{}
\[
\struck{\textstyle\prod} \ \operatorname{Hom}(X', I_{X'})^E \qquad\qquad \operatorname{Hom}(X', I_{X'}^E)
\]
\[
(P \amalg P)^E \qquad \struck{\ill{}}
\]
Topos associé à l'espace compact $I^E$
\note{la lettre notée $\mathbb{L}_X$ est un $\mathbb{L}$ ou un $\mathcal{L}$ cursif, non tranché}

\page{26}
\struck{Sous-espace algébrique fermé d'un \ill{} $=$ espace algébrique affine.}

\textbf{(9)} $\varprojlim$ dans la catégorie $\mathrm{Aff}_k$ des espaces algébriques affines sur $k$.

Correspond aux $\varinjlim$ dans la catégorie $\mathrm{Alg}_k$ des $k$-\emph{algèbres}.

\quad Bien adapté : l'interprétation de $C$ comme sous-catégorie pleine de $\widehat{C}$.

* Produits de schémas $\Longleftrightarrow$ « sommes » de $k$-algèbres $A_i$ ($i \in I$)
\[
A = \bigotimes_{i \in I} A_i
\]
Cas particulier : \struck{Objet final} $e_k = V_k$. \qquad \emph{Sections} sur $k$ (ou $V_k$, ou $e_k$) $=$ pts à valeurs dans $k$ \struck{\ill{}}

* Produit fibré $\Longleftrightarrow$ Produit tensoriel $B \otimes_A C$

\quad [produit dans $C_{/S}$] \qquad (« somme » dans $\mathrm{Alg}_A$)

$\Updownarrow$ Changement de base $C_{/S} \to C_{/S'}$, adjoint à droite du \emph{foncteur} « oubli de $S'$ au profit de $S$ »

\begin{tikzcd}
X \arrow[d] & X' \arrow[l] \arrow[d] & T \arrow[dl] \arrow[ll, bend right=30] \\
S & S' \arrow[l] &
\end{tikzcd}

\begin{tikzcd}
B \arrow[r] & B' \arrow[r, dashed] & C' \\
A \arrow[u] \arrow[r] & A' \arrow[u] \arrow[ur] &
\end{tikzcd}

\note{le $T$ est surchargé, peut-être biffé ; sa flèche courbe vers $X$ et la flèche pointillée $B' \dashrightarrow C'$ sont celles de la page}

\marginal{\emph{Interprétation en termes de foncteurs} : « restriction » du foncteur des $k$-algèbres à $k'$-algèbres (ou des $A$-algèbres à $A'$-algèbres)}

Cela donne toutes les $\varprojlim$. Autres exemples :

* Noyaux de doubles flèches \quad $X \underset{v}{\overset{u}{\rightrightarrows}} Y$
\[
\operatorname{Ker}(u, v) \simeq X \times_{(Y \times Y)} \Delta_Y
\]

\begin{tikzcd}
X \arrow[r, "{(u, v)}"] & Y \times Y \\
K \arrow[u] \arrow[r] & Y \arrow[u, "\Delta_Y"']
\end{tikzcd}

Conoyaux de doubles flèches \struck{\ill{}} \quad $B \underset{v}{\overset{u}{\leftleftarrows}} A$
\[
\simeq B / \text{idéal engendré par les } u(x) - v(x)
\]
[

\begin{tikzcd}
B \arrow[d] & A \otimes A \arrow[l, "{(u, v)}"'] \arrow[d] \\
C = ? & A \arrow[l]
\end{tikzcd}

$x \otimes y \mapsto xy$ : \quad morphisme surjectif ayant noyau engendré par $x \otimes 1 - 1 \otimes x$

* limites projectives filtrantes : correspond aux limites inductives filtrantes dans $\mathrm{Alg}_k$, qui se construisent ensemblistement.

\emph{Exemple} Noyau d'un hom. de groupes algébriques affines sur $k$
\[
G \longrightarrow G'
\]
\emph{Ex} \quad $\mathrm{Sl}(n)_k = \operatorname{Ker}\bigl(\mathrm{Gl}(n)_k \xrightarrow{\ \det\ } \mathbb{G}_{m,k}\bigr)$. \quad \struck{\ill{}}

\page{27}
$G$ opère sur $X$, \emph{stabilisateur} d'un $x \in X(k) \simeq \operatorname{Hom}_{\mathrm{Aff}_k}(e_k, X)$
\[
g \longmapsto g.x
\]

\begin{tikzcd}
G \arrow[r] & X \\
G_{(x)} \arrow[u] \arrow[r] & e_k \arrow[u]
\end{tikzcd}

\note{le coin inférieur gauche porte d'abord une autre lettre, noircie, puis $G_{(x)}$ récrit dessous}

$I$.

Si on veut le sous-groupe de stabilité de plusieurs éléments $x_i \in X_i(k)$, $G$ opérant sur les $X_i$, on prend $\bigcap_i G_{x_i}$ (cas particulier de produit fibré) où on fait opérer $G$ sur $X = \prod X_i$, et on prend le stabilisateur de $x = (x_i)_{i \in I} \in X(k)$.

\textbf{(10)} $\varinjlim$ dans $\mathrm{Aff}_k$ : Ne s'interprète pas de façon commode
en termes de $\widehat{\mathrm{Aff}_k} \simeq \underline{\operatorname{Hom}}(\mathrm{Alg}_k, (\mathrm{Ens}))$, et n'a pas beaucoup de sens géométrique, sauf en des cas très particuliers.

[\emph{Exemple} \quad $G = \mathrm{Gl}(n)_k$, $B = \mathrm{Tr}(n)_k \subset G$, alors $B$ définit une relation d'équivalence dans $G$, et le quotient $G/B$ \emph{dans} $\mathrm{Aff}_k$ est $e_k$ --- ce qui est idiot !]
\marginal{définir la relation d'équivalence \uncertain{définie par un} sous-groupe \ill{} d'un groupe \ill{}}

Correspond : $\varprojlim$ dans $\mathrm{Alg}_k$. Ce sont les $\varprojlim$ ensemblistes (nous fait une belle jambe !)

* \emph{Sommes dans} $\mathrm{Aff}_k$. Correspond aux produits dans
$\mathrm{Alg}_k$. Seules les sommes finies ($\Longleftrightarrow$ produits finis) ont une interprétation \emph{géométrique simple}.
\[
X_i \longrightarrow X \qquad (i \in I)
\]
Attention, $\coprod_i X_i(k') \to X(k')$ n'est pas bijective pour $\forall\, k'$, i.e.
\[
\coprod_i \operatorname{Hom}_{k\text{-alg}}(A_i, k') \longrightarrow \operatorname{Hom}_{k\text{-alg}}\Bigl(A = \prod A_i,\ k'\Bigr)
\]
\uncertain{est seulement} injective si $k' \neq 0$ \struck{\ill{}}, non pas bijective. (Un hom. de $A$ dans $k'$ « consiste » en la donnée d'une famille d'idempotents $e_i$ de $k'$ ($i \in I$), de somme 1 et deux à deux orthogonaux, et d'hom. $A_i \to k' / (1 - e_i) k'$ … Mais

\page{28}
$k'$ est « connexe » i.e. \uncertain{si on les} \uncertain{seules} décomposition de $1 \in k'$, (on a $I \neq \emptyset$ et) \uncertain{sont} les \uncertain{triviales}, l'une étant 1, \struck{\ill{}} (p. ex. $k'$ intègre, \struck{\ill{}} ou $k'/\mathrm{Nil}_{k'}$ intègre, \struck{\ill{}} ou \uncertain{si} $k'$ anneau local…) alors les applications envisagées sont bijectives.
\marginal{équivalent à la notion de connexité dans une catégorie… --- Exemple : $k'$ intègre}

Si $I = \emptyset$, on trouve un objet initial dans $\mathrm{Aff}_k$
\[
\emptyset_k = V_0 \quad (\text{algèbre nulle sur } k)
\]
\[
\emptyset_k(k') =
\begin{cases}
\emptyset & \text{si } k' \neq 0 \quad \text{i.e. } V_{k'} \not\simeq \emptyset_k \\
\{e\} & \text{si } k' = 0
\end{cases}
\qquad\qquad \emptyset_k(k') = \operatorname{Hom}(V_{k'}, \emptyset_k)
\]
\uncertain{On}. \struck{\ill{}}
\note{« $I \neq \emptyset$ » est ajouté au-dessus de la ligne, avec un renvoi ; « $\emptyset_k(k') = \operatorname{Hom}(V_{k'}, \emptyset_k)$ » est écrit sous la formule, relié par un signe $=$ vertical}

Donc c'est un \emph{objet initial strict}. On l'appelle le \emph{$k$-espace algébrique vide} (\struck{objet} objets de catégorie…,

(*) Conoyaux $\operatorname{Coker}(u, v : R \rightrightarrows X)$

\quad correspond à \quad $\operatorname{Ker}(u, v : A \rightrightarrows B)$

\quad mais n'a pas beaucoup de sens. O.K. si $R$ est une « relation d'équivalence » dans $X$ et si $R$ \uncertain{fini} sur $X$ par $u$, i.e. $B$ une $A$-algèbre finie. On reviendra dessus avec la question des passages au quotient.

Sous-espaces algébriques fermés d'un espace algébrique affine $\mathfrak{X} \simeq V_A$ : ils correspondent aux \emph{idéaux} de $A$, avec relation d'ordre renversée.

\textbf{(\uncertain{11})} \emph{Spectre d'une $k$-algèbre $A$} : espace sous-jacent.
\[
A \longleftrightarrow \mathfrak{X} = V_A
\]
On va associer un objet géométrique

(a) Ensemble sous-jacent : ensemble \struck{\ill{}} des « lieux » (loci) « de » pts géométriques de $\mathfrak{X}$ $\Longleftrightarrow$ ensemble \struck{\ill{}} $\operatorname{Spec}(A)$ des idéaux premiers de $A$
\note{le chiffre cerclé est surchargé (« 1 » suivi d'un second chiffre repassé) ; « $\operatorname{Spec}(A)$ » est ajouté au-dessus de la dernière ligne}

\page{29}
La connaissance de l'ensemble $X$ et de la famille des corps $(k(x))_{x \in X}$ (corps résiduels \struck{\ill{}} des pts de $\operatorname{Spec} A$ ou $\operatorname{Spec}(\mathfrak{X})$ --- \uncertain{k(x)} représente le foncteur $k'$ (corps variable sur $k$) $\longmapsto$ ens. des \struck{points} éléments de $\mathfrak{X}(k')$ \uncertain{associés} pour le lieu $k(x)$).

\textbf{(b)} \emph{Topologie sur $X$} (les fermés sont les) $V(S) \subset X$ associés aux parties $S = (f_\alpha)$ \struck{de} $A$, \struck{\ill{}} un point géom. \uncertain{à valeurs} dans $k'$ \uncertain{est} dans $V(S)$ \struck{\ill{}} si $f_\alpha(u) = 0$ pour tout $\alpha$. On a \struck{$V(S) = V(J_S)$ si}
\[
V(S) = \text{ensemble des idéaux premiers de } A \text{ qui contiennent } S.
\]
On a $V(S) = V(J_S)$, $J_S$ idéal engendré, et \uncertain{même} $V(S) = V(\widetilde{J}_S)$. \struck{Donc} Il est bon de se borner aux idéaux, et \uncertain{même} aux idéaux « radicaux ». \struck{L'ensemble des} On voit que $V(S) = V(S')$ (resp. $V(S) \subset V(S')$) signifie, vu des notations précédentes, que $V'_S = V'_{S'}$ (restrictions aux corps) \struck{ou encore que} (resp. $V'_S \subset V'_{S'}$), donc $\widetilde{J}_S = \widetilde{J}_{S'}$ (resp. $\widetilde{J}_S \supset \widetilde{J}_{S'}$) : car prouvons en fait que $\widetilde{J}_S$ est l'intersection des idéaux premiers qui contiennent $A$\note{sic : on attend « qui contiennent $S$ »}.
\[
\struck{V(\textstyle\sum J)} \quad V\Bigl(\bigcup_{i \in I} S_i\Bigr) = V\Bigl(\sum J_i\Bigr) = \bigcap_i V(S_i) = \bigcap_i V(J_i) \qquad I \text{ quelc.}
\]
(NB une \struck{intersection} somme de deux idéaux \struck{\ill{}} radicaux n'est pas nécessairement un idéal radical)
\[
V\Bigl(\bigcap_i J_i\Bigr) = V\Bigl(\prod J_i\Bigr) = \bigcup_i V(J_i) \qquad I \text{ fini}
\]
[NB $\prod_i J_i \subset \bigcap_i J_i$, il faut prouver que tt idéal premier qui contient \struck{l'intersection} (resp. produit) des $J_i$ contient un des $J_i$, on raisonne par l'absurde (i.e. si $\exists f_i \in J_i \setminus \mathfrak{p}$), \struck{on en forme \ill{} $A/\mathfrak{p}$ qui n'a pas de diviseurs de zéro} $\prod f_i$ \struck{\ill{}} $\in \prod J_i$ et $\notin \mathfrak{p}$, absurde] \struck{l'un d'eux est nul}

Donc on a une \emph{topologie}. Si $x \in X$ correspond à $\mathfrak{p}_x \subset A$, alors $\overline{\{x\}}$ est le plus petit fermé $V(J)$ qui contient $x$, i.e. tel que $J \subset \mathfrak{p}_x$, donc \struck{\ill{}} c'est $V(\mathfrak{p}_x)$.
\marginal{$x \in \overline{\{y\}} \Longleftrightarrow \overline{\{x\}} \subset \overline{\{y\}} \Longleftrightarrow \mathfrak{p}_x \supset \mathfrak{p}_y$}

\emph{Espace \struck{\ill{}} quasi-compact} \quad (\ill{})

[$X_f = X - V(f)$ est ouvert (le lieu des $u \in \mathfrak{X}(k)$ \ill{} dans $X_f$ est $f(u) \neq 0$), et les $X_f$ forment une base de la topologie de $X$, stable par intersections finies ($X_1 = X$, $X_{fg} = X_f \cap X_g$).
\marginal{plus bas}

\page{30}
[$X$ quasi-compact $\Longleftrightarrow$ si $A = \sum_{i \in I} J_i$, $\exists J$ fini $\subset I$, $A = \sum_{i \in J} J_i$.

\struck{\ill{}}

Si $J$ est un idéal de $A$, $\operatorname{Spec}(A/J) = Y \subset X$ est un sous-espace topologique fermé de $X$. \struck{On trouve \ill{}} $Y$ ne change pas si on remplace $J$ par $\widetilde{J}$, donc \struck{\ill{}} tout sous-espace fermé de $X$ est toujours homéo. au spectre d'un anneau réduit.

\emph{Prop} Si $A$ est réduit, $X = \operatorname{Spec} A$ irréductible $\Longleftrightarrow$ $A$ intègre.

En effet, $\operatorname{Spec} A$ irréductible
\[
\overset{\text{déf}}{\Longleftrightarrow}
\begin{cases}
X \neq \emptyset \\
X = X' \cup X'' \Rightarrow X' = X \text{ ou } X'' = X \quad (X', X'' \text{ fermés})
\end{cases}
\quad \text{i.e.} \quad
\begin{cases}
A \neq 0 \\
\text{i.e. } J' . J'' = 0 \Rightarrow J' \text{ ou } J'' \text{ nul}
\end{cases}
\]

\emph{Corollaire} Soit $A$ quelconque, \struck{$A$ a une correspondance biunivoque} les parties fermées irréductibles de $A$ sont les $V(\mathfrak{p})$, $\mathfrak{p}$ idéal premier de $A$ (correspondance biunivoque renversant l'ordre) et si $x$ correspond à $\mathfrak{p}$, alors $V(\mathfrak{p})$ est égal à $\overline{\{x\}}$. En particulier, $X$ est \emph{sobre}.

\quad \emph{NB} L'espace topologique $\operatorname{Spec} A$ ne dépend que de la structure d'anneau de $A$, i.e. de $\mathbb{Z}$-algèbre…

\emph{Spectre d'une $k$-algèbre : faisceau d'anneaux}

On ne peut reconstituer $A$ à partir du seul espace topologique $X = \operatorname{Spec} A$ !

Faisceaux sur $X$ : \quad $\forall f \in A$, ensemble $P(f)$

$\forall f, g \in A$ tels que $X_f \subset X_g$, une restriction $P(g) \to P(f)$
\[
\begin{aligned}
&\text{i.e. } V(f) \supset V(g) \\
&\text{i.e. } \widetilde{fA} \subset \widetilde{gA} \\
&\text{i.e. } f \in \widetilde{gA} \\
&\text{i.e. } \begin{cases} \exists n \in \mathbb{N}, \\ h \in A \end{cases} f^n = gh \qquad \Big| \text{ écrivons } f \prec g
\end{aligned}
\]
avec transitivité si $f, g, h$ tels que $f \prec g \prec h$ \quad et \uncertain{identification} si $f = g$

Avec $\forall f \in A$ et $(f_i)_{i \in I}$, $f_i \in I$ tels que
\[
\begin{aligned}
&X_f = \bigcup X_{f_i} \\
&\text{i.e. } V(f) = \bigcap V(f_i) \\
&\text{i.e. } \widetilde{fA} = \widetilde{\textstyle\sum f_i A} \\
&\text{i.e. } \begin{cases} \forall i \ f_i \prec f \\ \exists n \in \mathbb{N} \text{ et } (h_i)_{i \in I} \text{ avec } f^n = \sum f_i h_i \end{cases}
\end{aligned}
\]
on veut exactitude dans
\[
P(f) \longrightarrow \prod_i P(f_i) \rightrightarrows \prod_{ij} P(f_i f_j)
\]
\marginal{$\operatorname{Spec}(A_f) \simeq X_f$ (homéomorphisme) --- 10') Fonctorialité de \struck{\ill{}} $\operatorname{Spec} A$ \uncertain{tr.} d'anneaux --- 11)}
\note{« $f_i \in I$ » est sur la page, pour $i \in I$ ; la marge renvoie par une flèche à la ligne « $\operatorname{Spec} A_f \simeq X_f$ » de la page 21}

\page{31}
[\emph{NB} On peut supposer $I$ fini si on y tient]

Soit $M$ un $A$-module. On définit, pour $f \in A$
\[
M_f = \varinjlim_n (M, n) \qquad \text{où} \qquad (M, n) \xrightarrow{\ f^{n'-n}\ } (M, n') \quad \text{si } n' \geqslant n, \qquad (M,n) = M
\]
Donc les éléments de $M_f$ sont les $m/f^n$, avec règles de calcul habituelles. C'est la solution au pb universel d'envoyer $M$ dans un $A$-module où $f$ opère de façon inversible.
\[
M \xrightarrow{\ i_f\ } M_f
\]

\emph{Lemme} Si $f \prec g$, $\exists$ un unique hom. de $M_g$ dans $M_f$ compatible avec $i_g$, $i_f$

\begin{tikzcd}
& M \arrow[dl, "i_g"'] \arrow[d, "i_f"] \\
M_g \arrow[r, "i_{f,g}"'] & M_f
\end{tikzcd}

En effet, $f$ devenant inversible dans $M_f$ \uncertain{et} $f^n = gh$ aussi, $g$ devient inversible dans $M_f$, et ceci est l'assertion.

\emph{Cor} Transitivité : on a un préfaisceau $\widetilde{M}$ sur la base des \ill{}, catégorie \struck{rigide} \ill{} $A^\circ$. En fait, c'est un préfaisceau de modules sur le préfaisceau d'anneaux $f \mapsto A_f$.

\marginal{\emph{NB} $A_f$ une $A$-algèbre, $A_f \to A_g$ hom. d'algèbres, $M_f$ un $A_f$-module [$A_f$ \ill{} ; $M_f \simeq M \otimes_A A_f$] ; si $g \prec f$, $M_g = (M_f)_g$ \ill{} des transitivités}

\emph{\ill{}} Le préfaisceau précédent est un \emph{faisceau}.
\[
M_f \longrightarrow \prod M_{f_i} \rightrightarrows \prod M_{f_i f_j} \qquad \text{exact si } f_i \prec f \text{ et } f^n = \sum f_i h_i
\]
On est ramené au cas où $f = 1$ \struck{\ill{} $M_f$} (car $f \succ g$ implique $(M_f)_g \simeq M_g$)

…

En particulier, $X$ est \emph{annelé par} $\widetilde{A}$, noté $\underline{\mathcal{O}}_X$, faisceau \struck{\ill{}} d'anneaux (et \uncertain{même} de $A$-algèbres).

Fibres du faisceau \quad $M_{\mathfrak{p}}$, $A_{\mathfrak{p}} =$ \quad [$S \subset A$ stable par multiplication donc filtrant pour \struck{\ill{}} $\succ$, donc $MS^{-1} = \varinjlim_{f \in S} M_f = \{ \frac{x}{f^n} \ldots$ règles de calcul…]
\[
A_{\mathfrak{p}} = \underline{\mathcal{O}}_{X,x} \quad \text{anneau local}
\]

\end{document}
